\subsubsection*{Selección de Artículos}

El proceso de selección se ejecutó conforme al flujo PRISMA 2020, articulando de forma secuencial las fases de identificación, cribado, obtención de texto completo, elegibilidad e inclusión final. Los criterios operacionales de inclusión y exclusión utilizados en las fases de cribado y elegibilidad se presentan en la Tabla~\ref{tab:criterios_ie}. Tras la consolidación de resultados provenientes de las bases consultadas y la depuración de duplicados, se conformó el conjunto de registros a cribar por título y resumen.

El cribado de título y resumen se aplicó de manera sistemática sobre \texttt{articles\_screened.csv}, evaluando para cada registro la relación con el fenómeno de calidad del aire (P), el uso central de lógica difusa (I), la evaluación de calidad del aire mediante lógica difusa (P+I), la naturaleza del estudio y la existencia de plataforma o sistema de monitoreo (Co). Esta fase permitió separar los estudios excluidos de aquellos que debían pasar a recuperación de texto completo.

La obtención de textos completos combinó una estrategia automatizada y una estrategia manual. En la fase automatizada, cada registro con DOI se consultó primero en Unpaywall y, cuando no se identificó un PDF directo, se aplicó una segunda consulta en OpenAlex para localizar enlaces de acceso abierto. Esta etapa permitió clasificar de forma trazable los casos descargados y los no recuperables por ausencia de DOI, falta de enlace PDF directo o restricción de acceso. Para los casos no obtenidos automáticamente, se realizó búsqueda manual en las plataformas de las bases indexadas, en arXiv y en ResearchGate; cuando el documento no estuvo disponible en acceso directo, se gestionó la solicitud a los autores mediante ResearchGate. El estado de recuperación se documentó en \texttt{articles\_screened.csv} a través de \texttt{Retrieval\_result} y en la bitácora operativa de descargas (\texttt{downloads.csv}).

La elegibilidad se evaluó sobre texto completo en \texttt{articles\_full\_text.csv}, aplicando los criterios D1--D3 y el criterio transversal de diseño metodológico y validación, con registro explícito de evidencia por página. En esta etapa se consignó una razón única de exclusión por estudio cuando correspondía, manteniendo consistencia con el reporte PRISMA de exclusiones en texto completo. El detalle cuantitativo por fase se presenta en la sección de resultados y en el diagrama PRISMA, preservando aquí el foco metodológico del procedimiento.
